\documentclass[10pt]{article}
\usepackage[T1]{fontenc}
\usepackage[utf8]{inputenc}
\usepackage[french]{babel}
\usepackage{amsmath,amssymb,mathrsfs,mathpazo}
\usepackage{geometry}
\geometry{a4paper,margin=2cm}
\usepackage{graphicx,array,tabularx,tabularray,colortbl,multirow,multicol,booktabs,enumitem,pifont,ulem,fancybox,diagbox,cancel}
\usepackage{pstricks,pstricks-add,pst-plot,pst-tree,pst-node,pst-text,pst-func,pst-math,pst-bspline,pst-3dplot,pst-eucl,pst-solides3d}
\let\pscalculate\undefined
\usepackage{pst-bezier}
\usepackage{tikz,pgfplots,tkz-tab}
\usetikzlibrary{arrows,arrows.meta,patterns,calc,positioning,trees,decorations.markings,shapes.geometric,shapes.misc,shapes.arrows,intersections,decorations.pathmorphing,decorations.pathreplacing,backgrounds,shadows,babel}
\pgfplotsset{compat=1.18}
\usepgfplotslibrary{fillbetween}
\usepackage[np]{numprint}
\usepackage{listings,makecell,pas-tableur,twoopt}
\let\barre\undefined
\usepackage[pstricks]{tablvar}
\definecolor{darkgreen}{RGB}{0,100,0}
\providecommand{\red}{\color{red}}
\providecommand{\blue}{\color{blue}}
\providecommand{\green}{\color{green}}
\providecommand{\black}{\color{black}}
\providecommand{\white}{\color{white}}
\providecommand{\gray}{\color{gray}}
\providecommand{\cyan}{\color{cyan}}
\providecommand{\magenta}{\color{magenta}}
\providecommand{\yellow}{\color{yellow}}
\providecommand{\violet}{\color{violet}}
\providecommand{\orange}{\color{orange}}
\providecommand{\euro}{\texteuro}
\providecommand{\EUR}{\texteuro}
\providecommand{\ieme}{\textsuperscript{e}}
\providecommand{\degres}{\ensuremath{^\circ}}
\providecommand{\degree}{\ensuremath{^\circ}}
\providecommand{\vv}[1]{\overrightarrow{#1}}
\providecommand{\Coord}[1]{#1}
\providecommand{\decosix}{}
\providecommand{\starredbullet}{\ensuremath{\bullet}}
\pagestyle{empty}
\setlength{\parindent}{0pt}
\usepackage{adjustbox}\DeclareUnicodeCharacter{2236}{:}\DeclareUnicodeCharacter{2190}{\ensuremath{\leftarrow}}
\providecommand{\R}{\mathbb{R}}
\providecommand{\N}{\mathbb{N}}
\providecommand{\Z}{\mathbb{Z}}
\providecommand{\Q}{\mathbb{Q}}
\providecommand{\D}{\mathbb{D}}
\providecommand{\C}{\mathbb{C}}
\providecommand{\e}{{\mathrm{e}}}
\providecommand{\vect}[1]{\overrightarrow{#1}}
\providecommand{\vectt}[1]{\overrightarrow{\mathrm{#1}}}
\providecommand{\vvt}[1]{\overrightarrow{\mathrm{#1}}}
\providecommand{\np}[2][]{{#2~#1}}
\providecommand{\vv}[1]{{\overrightarrow{#1}}}
\renewcommand{\d}{{\mathrm{d}}}
\providecommand{\degres}{{^\circ}}
\providecommand{\degre}{{^\circ}}
\providecommand{\blue}{\color{blue}}
\providecommand{\red}{\color{red}}
\providecommand{\enskip}{\quad}
\providecommand{\llbracket}{\lbrack\!\lbrack}
\providecommand{\rrbracket}{\rbrack\!\rbrack}
\providecommand{\euro}{{\text{€}}}
\providecommand{\no}{{\text{n°}}}
\providecommand{\interval}[2]{{[#1;#2]}}
\providecommand{\intervFO}[2]{{[#1;#2[}}
\providecommand{\dotfill}{\ldots}

\providecommand{\ds}{\displaystyle}
\providecommand{\textding}[1]{#1}
\providecommand{\Oij}{\(\mathrm{O};\vec{\imath},\vec{\jmath}\)}
\providecommand{\Oijk}{\(\mathrm{O};\vec{\imath},\vec{\jmath},\vec{k}\)}
\providecommand{\Ouv}{\(\mathrm{O};\vec{u},\vec{v}\)}
\providecommand{\intervFF}[2]{\ensuremath{[#1;#2]}}
\renewcommand{\euro}{\ifmmode\text{\texteuro}\else\texteuro\fi}
\DeclareUnicodeCharacter{00B0}{\ifmmode^{\circ}\else\textdegree\fi}
\DeclareUnicodeCharacter{20AC}{\ifmmode\text{\texteuro}\else\texteuro\fi}
% Énoncés : Éducation nationale. Transcription et corrigés : APMEP, Denis Vergès et contributeurs. Découpage et classement : CoopMaths, Sébastien Lozano, Éric Elter et contributeurs. Adaptation de la présentation : Maths Explorer.
% CC BY-SA : https://coopmaths.fr/www/about/
% https://coopmaths.fr/alea/static/eam/2026/tex/eam_specialite_2026_04_polynesie_3.tex
\begin{document}
\noindent\textit{APMEP / CoopMaths. CC BY-SA. Présentation adaptée par Maths Explorer.}\par\medskip
Les activités d'une association sont financées grâce à des donateurs.

On réalise une estimation du nombre de donateurs de la manière suivante.

Chaque année,
\begin{itemize}[left=5mm]
	\item 80\,\% des donateurs de l'année précédente renouvellent leur don;
	\item 300 nouveaux donateurs effectuent un don.
\end{itemize}

En 2025, l'association comptait \np{1000} donateurs.

Pour tout entier naturel $n$, on modélise par $u_{n}$ le nombre de donateurs pour l'année $2025+n$.

Ainsi, on a $u_{0}=1000$.

\subsection*{Partie A}

\begin{enumerate}
	\item \begin{enumerate}
		\item Déterminer le nombre de donateurs prévu en 2026.

		\item Justifier que pour tout entier naturel $n$, on a :\quad $u_{n+1}=0,8 u_{n}+300$.
	\end{enumerate}

	\item On donne $u_{5} \approx \np{1336}$. Que représente cette valeur dans le contexte de l'exercice ?

	\item On considère la suite d'instructions Python ci-dessous :

\begin{lstlisting}[style=python,xleftmargin=5cm,xrightmargin=5cm,numbers=none]
u=1000
i=0
while u<=1400:
	u=0.8*u + 300
	i=i+1
\end{lstlisting}

	À la fin de l'exécution de ce script, la variable i contient la valeur 8.\\
	Comment peut-on l'interpréter dans le cadre de l'exercice?
\end{enumerate}

\subsection*{Partie B}

\begin{minipage}{11cm}
	On considère la suite $\left(v_{n}\right)$ définie sur $\N$ par \quad $v_{n}=u_{n}-\np{1500}$.

\begin{enumerate}
	\item Démontrer que $\left(v_{n}\right)$ est une suite géométrique.

	Préciser sa raison et son premier terme.
\end{enumerate}
\end{minipage}
\hfill
\begin{tblr}{colspec={X[c,4cm]},hlines, vlines}
	Aide au calcul :
	\[\dfrac{1200}{0,8}=\np{1500}\]
	\[0,8 \times \np{1500}=\np{1200}\]
\end{tblr}

\begin{enumerate}[start=2]
	\item \begin{enumerate}
		\item Exprimer $v_{n}$ en fonction de $n$.

		\item En déduire, pour tout entier naturel $n$, que : \quad $u_{n}=\np{1500}-500 \times 0,8^{n}$.
	\end{enumerate}
\end{enumerate}

\end{document}
